\documentclass[10pt,a4paper]{amsart}
\usepackage[margin=19mm]{geometry}
\usepackage{amsmath,amssymb,mathtools}
\usepackage[scr=rsfs,cal=euler]{mathalfa}
\usepackage{xfrac}
\usepackage[unicode,hidelinks,bookmarks=false]{hyperref}
\newcommand{\ie}{i.e.\ }
\newcommand{\cf}{cf.\ }
\newcommand{\resp}{respectively\ }
\newcommand{\Subsection}{\subsection}
\providecommand{\nobreakdash}{\nobreak\hbox{-}}
\setlength{\emergencystretch}{2em}
\multlinegap0pt
\usepackage{amssymb}
\usepackage{tikz}
\newtheorem{lemma}{Lemma}
\title{Time complexity for deterministic string machines}
\author{Nur \c{C}ataltepe, Vanessa Kosoy}
\date{Source: arXiv:2405.06043v2 (CC BY 4.0)}
\begin{document}
\maketitle

\noindent\emph{Source and licence.} Time complexity for deterministic string machines. Version arXiv:2405.06043v2 (CC BY 4.0), distributed by arXiv under the Creative Commons Attribution 4.0 International licence, http://creativecommons.org/licenses/by/4.0/, as recorded in the arXiv licence metadata for this exact version and shown on the abstract page https://arxiv.org/abs/2405.06043v2. Full licence notice: \texttt{licenses/arxiv-2405.06043-CC-BY-4.0.txt}.\par\smallskip

\noindent\textit{Editorial scope.} Counted unit: the article's lemma lem:fsmprepending (source0.tex lines 1074--1104). The article's complete original proof of it is delivered with it (source0.tex lines 1105--1165, one \texttt{proof} environment of 61 lines). No mathematics is written, repaired or re-derived: the statement and the proof are the article's own lines, in the article's own notation, and no retained line is substituted or re-flowed.  No further source-local context is needed: every cross-reference of every retained line resolves inside the two delivered spans.  Nothing was omitted from any retained span, so the package carries no omission record.  No retained line contains a citation, so the wrapper carries no bibliography and no bibliography entry.  No retained line contains a figure environment, an image-inclusion command or drawing code, so no image is delivered and the package declares no asset dependency.  Editorial formatting only, in addition: an AMS \texttt{amsart} preamble that restates the article's own declarations and macros used by the retained lines, namely the environments \texttt{lemma} on separate counters, together with the standard packages those lines need.  The retained environments print in this document as Lemma 1 in order of appearance, whereas the article prints the same items with the numbering of its own full document.  The complete native source of the article as distributed in the arXiv e-print for this exact version is bundled unchanged as \texttt{sources/158928/source0.tex}.
\begin{lemma}
    \label{lem:fsmprepending} Suppose some finite-state transducer
    $T$ with
    primary input signature
    $X\to X$ in a tape category $\mathcal{X}$, structure
    functor $F$, and output state space
    $S(A)$, such that these are its only inputs and outputs.
    Fix a finite set of morphisms $\{g_1,\dots,g_n:X\to X\}$ in
    the freely-generated $\mathbb{N}^2$-filtered category over
    $\mathcal{X}$ with a single generating signature
    $X\to X$ of degree $x$, such that the linear coefficient of
    the degree of each $g_i$ is at most $1$ (that is, the generator
    corresponding to $X\to X$ appears at most once).
    Suppose
    we are able to construct a
    primary input morphism $g:X\to X$
    for this transducer in the following
    way:
    \begin{itemize}
        \item Our starting value for $g$ is $\mathrm{id}_X$.
        \item We can replace $g$ with
        with one of the previously-given $g_i$,
        with $g$ substituted for the
        generator morphism. We can perform this step a finite
        of times.
    \end{itemize}
    Then, given some state transition reached as a result of some
    such $g$,
    the amount of memory required to compute the output state of
    the transducer is $O(1)$ in the degree of $g$.
\end{lemma}

\begin{proof}
We will assume that one ``computational step'' involves performing one
of the operations involved in constructing $g$ described above. We
will demonstrate that computing the next state after
any of the operations performed requires storing an amount of information
that is $O(1)$ in the degree of $g$.
\begin{itemize}
\item If one of the $g_i$ does not have the generator morphism appear
at all, then $S(A)$ after this $g_i$ has replaced $g$ depends solely
on the $g_i$ used, and since the set of such $g_i$ is finite, it takes
a finite amount of space to store the state reached from the starting
state when one of these $g_i$ are evaluated. 
\item If one of the $g_i$ only involves post-composing an endomorphism
of $X$ into the generator, then storing the state reached from a given
state after any of these endomorphisms is post-composed takes a fixed
amount of memory, since the set of $g_i$ is finite.
\item If one of the $g_i$ involves only pre-composing and post-composing
endomorphisms of $T$ into $g$, then computing the resultant state after
the action of one of these $g_i$ only requires us to store $|S(A)|$
extra ``possible current states.'' That is, if we know what state we
would be in after the evaluation of $g$ had we started at any other
one of the $|S(A)|$ states of the transducer, then after computing
which new starting state we would have started at (which, since the
set of endomorphisms we pre-compose is finite, is equivalent to storing
a fixed-size transition table), we can then pick
the state that the copy of the transducer that started at that state
would be at after the action of $g$. This is therefore equivalent
to running $|S(A)|$ copies of a finite state automaton, where each
copy starts at a different state .
\item Suppose one of the $g_i$ is obtained by taking a morphism
$\alpha:X\to X^n$, post-composing morphisms into each branch (such that
one branch includes the generator), and then finally post-composing
a morphism $\beta:X^n \to X$. Even if
multiple ``branchings'' occur, since the generator where $g$ will
be substituted appears only once, we can still decompose this
into a morphism obtained by
first taking the monoidal product of $g$ with a number of fixed
morphisms such that the overall signature of their product is
$X^n \to X^n$, and then composing in $\alpha$ and $\beta$.
Now, we only need
$O(1)$ space to store the functions $S(F(X)) \to S(F(X^n))$ and
$S(F(X^n)) \to S(F(X))$ induced by $\alpha$ and $\beta$, since these
are functions between finite sets. Additionally, for the function
$S(F(X^n)) \to S(F(X^n))$ induced by the monoid product of $g$
with with a fixed morphism
$\gamma$ of signature $X^{n-1} \to X^{n-1}$, by
the property discussed at the start of this section, we only need to
be keeping track of the action of $F(g \otimes \mathrm{id}_{X^{n-1}})$
on $S(F(X^n))$ (which, since there are a finite number of
$g_i$, only requires keeping an absolutely bounded number of automata
running), and then compute the action of
$F(\mathrm{id}_X\otimes \gamma)$ on the result, which, as a function
with a finite domain and codomain, takes finite space to store.
\end{itemize}
Therefore, to compute the state transition induced by $g$, we only
need to keep $|S(A)|$ copies of the finite state automaton running,
as well as some fixed number of
finite-state automata with input and output
state spaces corresponding to
$S(F(X^n))$ for some fixed $n$.
\end{proof}

\end{document}
